\documentclass[10pt,a4paper]{amsart}
\usepackage[margin=19mm]{geometry}
\usepackage{amsmath,amssymb,mathtools}
\usepackage[scr=rsfs,cal=euler]{mathalfa}
\usepackage{xfrac}
\usepackage[unicode,hidelinks,bookmarks=false]{hyperref}
\newcommand{\ie}{i.e.\ }
\newcommand{\cf}{cf.\ }
\newcommand{\resp}{respectively\ }
\newcommand{\Subsection}{\subsection}
\providecommand{\nobreakdash}{\nobreak\hbox{-}}
\setlength{\emergencystretch}{2em}
\multlinegap0pt
\usepackage[shortlabels]{enumitem}
\usepackage{xcolor}
\usepackage{float}
\usepackage{amssymb}
\usepackage{mathtools}
\newtheorem{lemma}{Lemma}
\newcommand{\Z}{\mathbb Z}
\renewcommand{\ge}{\geqslant}
\renewcommand{\le}{\leqslant}
\title{Wold-type decompositions for norm-increasing $m$-concave weighted shifts on directed trees}
\author{Ashish Kujur}
\date{Source: arXiv:2609.14662v1 (CC BY 4.0)}
\begin{document}
\maketitle

\noindent\emph{Source and licence.} Wold-type decompositions for norm-increasing $m$-concave weighted shifts on directed trees. Version arXiv:2609.14662v1 (CC BY 4.0), distributed by arXiv under the Creative Commons Attribution 4.0 International licence, http://creativecommons.org/licenses/by/4.0/, as recorded in the arXiv licence metadata for this exact version and shown on the abstract page https://arxiv.org/abs/2609.14662v1. Full licence notice: \texttt{licenses/arxiv-2609.14662-CC-BY-4.0.txt}.\par\smallskip

\noindent\textit{Editorial scope.} Counted unit: the article's lemma lem:Cesaro (source0.tex lines 525--531). The article's complete original proof of it is delivered with it (source0.tex lines 532--562, one \texttt{proof} environment of 31 lines). No mathematics is written, repaired or re-derived: the statement and the proof are the article's own lines, in the article's own notation, and no retained line is substituted or re-flowed.  No further source-local context is needed: every cross-reference of every retained line resolves inside the two delivered spans.  Nothing was omitted from any retained span, so the package carries no omission record.  No retained line contains a citation, so the wrapper carries no bibliography and no bibliography entry.  No retained line contains a figure environment, an image-inclusion command or drawing code, so no image is delivered and the package declares no asset dependency.  Editorial formatting only, in addition: an AMS \texttt{amsart} preamble that restates the article's own declarations and macros used by the retained lines, namely the environments \texttt{lemma} on separate counters, together with the standard packages those lines need.  The retained environments print in this document as Lemma 1 in order of appearance, whereas the article prints the same items with the numbering of its own full document.  The complete native source of the article as distributed in the arXiv e-print for this exact version is bundled unchanged as \texttt{sources/210344/source0.tex}.
\begin{lemma}
    Let $\left\{ a_n \right\}_{n \in \Z_{\ge 1}}$ be a nonnegative nondecreasing sequence which is not identically zero. Define $\sigma_n := \frac{1}{n} \sum_{k=1}^{n} a_k$ for each $n \in \Z_{\ge 1}$. Then we have
    \begin{equation*}
        \liminf_{n \to \infty} \frac{a_n}{\sigma_n ^{2}} < \infty.
    \end{equation*}
    \label{lem:Cesaro}
\end{lemma}

\begin{proof}
    First, we observe that since $\{ a_n \}_{n \in \Z_{\ge1}}$ is nondecreasing, $\{ \sigma_{n} \}_{n \in \Z_{\ge 1}}$ is also nondecreasing. Indeed, by monotonicity of $\{ a_n \}_{n \in \Z_{\ge1}}$, we have $\sigma_n \le a_{n+1}$ for each $n \in \Z_{\ge 1}$, and consequently,
    \begin{equation*}
        \begin{aligned}
            \sigma_{n+1} - \sigma_n = \frac{1}{n+1} \sum_{k=1}^{n+1} a_k - \sigma_n 
                                    = \frac{n \sigma_n + a_{n+1}}{n+1} - \sigma_n 
                                    = \frac{a_{n+1} - \sigma_n}{n+1} \ge 0,
        \end{aligned}
    \end{equation*}
    for all $n \in \Z_{\ge 1}$.

    Moreover, since $\left\{ a_n \right\}_{n \in \Z_{\ge 1}}$ is not identically zero we have that there is some $N \in \Z_{\ge 1}$ such that $\sigma_n > 0$ for each $n \in \Z_{\ge N}$.

    Now, observe that $a_n = n \sigma_n - (n-1) \sigma_{n-1} = \sigma_n + (n-1) (\sigma_n - \sigma_{n-1})$ for each $n \in \Z_{\ge 2}$. Using the fact that $\sigma_{n-1} \le \sigma_{n}$ for each $n \in \Z_{\ge 2}$, it follows that for each $n \in \Z_{\ge N+1}$,
    \begin{equation}
        \begin{aligned}
        \frac{a_n}{\sigma_n ^2} &= \frac{\sigma_n + (n-1) (\sigma_n - \sigma_{n-1})}{\sigma_n ^2} \\
        &= \frac{1}{\sigma_n} + \frac{(n-1) (\sigma_n - \sigma_{n-1})}{\sigma_{n} ^2} \\
        &\le \frac{1}{\sigma_n} + (n-1) \left( \frac{1}{\sigma_{n-1}} - \frac{1}{\sigma_n} \right).
    \end{aligned}
    \label{eq:Cesaro-estimate}
    \end{equation}
    
    We define $c_n := \frac{1}{\sigma_{n-1}} - \frac{1}{\sigma_n}$ for each $n \in \Z_{\ge N+1}$. Observe that $c_n \ge 0$ for each $n \in \Z_{\ge N+1}$ and $\sum_{n = N+1}^{\infty} c_n < \infty$. Consequently, $\liminf_{n \to \infty} nc_n = 0$. Indeed, otherwise there would exist some $\varepsilon > 0$ and some $N_0 \in \Z_{\ge N+1}$ such that $c_n \ge \varepsilon/n$ for all $n \in \Z_{\ge N_0}$, but this would imply that $\sum_{n = N+1}^{\infty} c_n = \infty$.

    We now select a strictly increasing sequence $\{ n_j \}_{j \in \Z_{\ge 1}} \subset \Z_{\ge N+1}$ such that $n_j c_{n_j} \to 0$. From \eqref{eq:Cesaro-estimate}, we have that 
    \begin{equation*}
        \frac{a_{n_j}}{\sigma_{n_j} ^2} \le \frac{1}{\sigma_N} + (n_j - 1) c_{n_j}, \qquad j \in \Z_{\ge 1}.
    \end{equation*}
    Since the right hand side of the previous inequality is bounded, we have $\liminf_{n \to \infty} \frac{a_{n}}{\sigma_{n} ^2} < \infty$. This completes the proof of the lemma.
\end{proof}

\end{document}
